\chapter*{Résumé Exécutif}
\addcontentsline{toc}{chapter}{Résumé Exécutif}

Le présent document constitue le \textbf{Guide de Préparation Intensive et Haute Rentabilité} destiné aux candidats au concours de recrutement des \textbf{Ingénieurs d'État de 1\textsuperscript{er} grade} du \textbf{Ministère de l'Économie et des Finances (MEF)} du Royaume du Maroc.

Face à la contrainte temporelle extrême (dernière ligne droite avant l'épreuve écrite du \textbf{20 septembre}), l'ensemble des contenus a été sélectionné et structuré selon le principe d'efficacité maximale : \textit{Recherche $\rightarrow$ Vérification $\rightarrow$ Extraction $\rightarrow$ Priorisation $\rightarrow$ Mémorisation ciblée}.

Ce guide s'articule autour de cinq blocs fondamentaux :
\begin{enumerate}[label=\textbf{\arabic*.}]
    \item \textbf{Cadrage et Méthodologie de l'Épreuve Écrite :} Analyse du format de l'examen (sujet de réflexion/dissertation et QCM), grille d'évaluation des correcteurs du MEF, gestion des trois heures, et méthode standardisée de construction de plan en deux parties équilibrées ($I.$ Diagnostic \& Enjeux, $II.$ Stratégies \& Perspectives).
    \item \textbf{Annales Réelles Décortiquées :} Analyse exhaustive des sujets réellement proposés lors des dernières sessions (session du 28 avril 2024 portant sur \textit{Maroc Digital 2030} et la \textit{Transition Énergétique} ; sessions précédentes portant sur \textit{l'Inflation}, le \textit{Déficit Budgétaire}, la \textit{Protection Sociale} et la \textit{Charte de l'Investissement}). Chaque sujet est décomposé selon la formule : \textit{Sujet $\rightarrow$ Problématique $\rightarrow$ Plan détaillé $\rightarrow$ Arguments clés $\rightarrow$ Exemples et chiffres marocains $\rightarrow$ Conclusion}.
    \item \textbf{Banque Exhaustive de Questions / Réponses Normalisée :} Recueil structuré de questions récurrentes issues des annales et concours d'ingénieurs et administrateurs du MEF, classées selon un code de priorité strict (🔴 Très important, 🟠 Important, 🟡 À connaître, ⚪ Secondaire) couvrant l'économie, les finances publiques (LOF 130-13), la modernisation administrative, l'organisation interne du MEF, et les systèmes d'information étatiques.
    \item \textbf{Fiches Mémotechniques « Haute Rentabilité » :} Synthèses condensées sur les missions des directions du MEF (DGI, TGR, Douanes, Budget, DTFE, DEPP, DEPF), les chiffres macroéconomiques clés, les grands principes budgétaires, et les orientations stratégiques du pays.
    \item \textbf{Plan d'Action Sprint Final :} Stratégie de travail par blocs horaires jusqu'au jour de l'examen pour maximiser la rétention et la performance rédactionnelle.
\end{enumerate}
